\documentclass[./main.tex]{subfiles}
\begin{document}

% Slide # 9
\begin{frame}[t, label=slide09]
        % Title
        \frametitle{Back to game theory}

        % Body
        \footnotesize

        Let us rewrite 

        \begin{align*}
                f(x) =&\,x(1-x)((a + d - b - c)x + (b-d)) = \nonumber \\
                     =&\,x(1-x)((a-c)x-(d-b)(1-x)) = \nonumber \\
                     =&\,x(1-x)(\alpha x - \beta(1-x)) \,.
        \end{align*}

        \vspace{3mm}

        \uncover<2->{
                The equilibria are thus $x_{1}=0$, $x_{2}=1$ and $x^{*}=\frac{\beta}{\beta+\alpha}$.
        }

        \vspace{5mm}

        \uncover<3->{
                We want to discard spurious solutions, therefore

                \begin{equation*}
                        x_{1}\leq x^{*}\leq x_{2}\;\Rightarrow\;\Gamma=\{(\alpha,\beta)\in \mathbb{R}^{2}\;:\;\text{sign}(\alpha)=\text{sign}(\beta)\}\,.
                \end{equation*}
        }
         
        \uncover<4->{
                By choosing $(\alpha,\beta)\in\Gamma$ carefully we can steer our population to adopt the desired strategy (equilibrium selection). 
        }

\end{frame}

\end{document}
